This observation is evidently quite general: formulas analogous to (57.2)
establish a correspondence between the eigenfunctions $\psi_{jm}$ and the
components of a covariant symmetric spinor of rank $2j$\footnote{The same
result can also be reached by a somewhat different argument.  Expand the
wave function $\psi$ of a particle whose angular momentum is $j$ in the
eigenfunctions $\psi_{jm}$:
$\psi=\sum_m a_m\psi_{jm}$.  The coefficients $a_m$ are then the probability
amplitudes associated with the different values of $m$.  In this sense, they
correspond to the ``components'' $\psi(m)$ of the spin wave function, and
this correspondence fixes their transformation law.  On the other hand, the
value of $\psi$ at a specified point in space cannot depend on the coordinate
system chosen; hence the sum $\sum a_m\psi_{jm}$ must be a scalar.  Comparison
with the scalar in (57.3) shows that the $a_m$ must transform in the same way
as $(-1)^{j-m}\psi_{j,-m}$.}:
\begin{equation}
 \psi_{jm}=\left[\frac{(2j)!}{(j+m)!(j-m)!}\right]^{1/2}
 \psi_{\overbrace{11\ldots}^{j+m}\overbrace{22\ldots}^{j-m}}.
 \tag{57.6}
\end{equation}
For integral angular momentum $j$, the eigenfunctions are the spherical
harmonics $Y_{jm}$.  The case $j=1$ is particularly important.
